\begingroup
\renewcommand{\arraystretch}{1.45}
\begingroup
\Needspace{45\baselineskip}
\InterpolationArticle{31}
Da die oben erläuterten Regeln eine Funktion $X'$ der Form
\[
\begin{aligned}
&\gamma+\gamma'\cos x+\gamma''\cos2x+.\,.\,.+\gamma^{m}\cos mx\\[5pt]
&\quad+\delta'\sin x+\delta''\sin2x+.\,.\,.\,.+\delta^{m}\sin mx
\end{aligned}
\]
liefern, die alle $\mu$ gegebenen Werte erfüllt und in der $m=\frac12\mu-\frac12$ oder $\frac12\mu$ gilt, je nachdem, ob $\mu$ ungerade oder gerade ist, lohnt es sich, diese Funktion mit der eben gefundenen Funktion
\[
\begin{aligned}
&\varepsilon+\varepsilon'\cos x+\varepsilon''\cos2x+\varepsilon'''\cos3x+.\,.\,.\\[5pt]
&\qquad{}+\varepsilon^{\mu-1}\cos(\mu-1)x
\end{aligned}
\]
\endgroup
\SourceLeaf{322}
zu vergleichen, die wir mit $X''$ bezeichnen werden. Vergleicht man die in den Artikeln 20, 22 angegebenen Werte der Koeffizienten
$\gamma,\gamma',\delta',\gamma'',\delta''$ usw. mit den eben gefundenen Werten der Koeffizienten
$\varepsilon,\varepsilon',\varepsilon''$ usw., so findet man
\[
\begin{aligned}
&\varepsilon=\gamma\\[5pt]
&\varepsilon'=\gamma'-\delta'\operatorname{cotang}\mu a,
\\[5pt]
&\varepsilon^{\mu-1}=\delta'\operatorname{cosec}\mu a\\[5pt]
&\varepsilon''=\gamma''-\delta''\operatorname{cotang}\mu a,
\\[5pt]
&\varepsilon^{\mu-2}=\delta''\operatorname{cosec}\mu a\\[5pt]
&\varepsilon'''=\gamma'''-\delta'''\operatorname{cotang}\mu a,
\\[5pt]
&\varepsilon^{\mu-3}=\delta'''\operatorname{cosec}\mu a\\[5pt]
&\text{usw., bis zu}\\[5pt]
&\left(\tfrac12\right)\varepsilon^m=\gamma^m-\delta^m\operatorname{cotang}\mu a,
\\[5pt]
&\left(\tfrac12\right)\varepsilon^{\mu-m}=\delta^m\operatorname{cosec}\mu a
\end{aligned}
\]
\Needspace{20\baselineskip}
Der Faktor $\tfrac12$, der $\varepsilon^m$ und $\varepsilon^{\mu-m}$ vorangestellt ist, gilt nur, wenn $\mu$ gerade ist, und muss weggelassen werden, wenn $\mu$ ungerade ist. Auf diese Weise sind im letzteren Fall, in dem $\mu-m=m+1$ gilt, offenbar alle $\mu$ Koeffizienten $\varepsilon,\varepsilon',\varepsilon''$ usw. durch $\gamma,\gamma',\delta'$ usw. bestimmt; im ersteren dagegen, in dem $m=\mu-m$ gilt, haben wir für den Koeffizienten $\varepsilon^{\frac12\mu}$ zwei Werte $2\gamma^{\frac12\mu}-2\delta^{\frac12\mu}\cot\mu a$, $2\delta^{\frac12\mu}\operatorname{cosec}\mu a$, deren Gleichheit sich aus der Gleichung $\gamma^{\frac12\mu}\sin\tfrac12\mu a=\delta^{\frac12\mu}\cos\tfrac12\mu a$ leicht erkennen lässt und für die wir daher auch diesen Wert $\gamma^{\frac12\mu}-\delta^{\frac12\mu}\operatorname{cotg}\mu a+\delta^{\frac12\mu}\operatorname{cosec}\mu a$ annehmen können.

\Needspace{25\baselineskip}
Hieraus lässt sich leicht ableiten, dass in beiden Fällen
\[
\begin{aligned}
&X''-(X'-\delta'\sin x-\delta''\sin2x-\delta^m\sin3x-\text{usw.}-\delta^m\sin mx)\\[5pt]
&=-\operatorname{cotg}\mu a(\delta'\cos x+\delta''\cos2x+\delta'''\cos3x+\mathinner{.\,.\,.}+\delta^m\cos mx)\\[5pt]
&\quad+\operatorname{cosec}\mu a(\delta'\cos(\mu-1)x+\delta''\cos(\mu-2)x+\mathinner{.\,.\,.}+\delta^m\cos(\mu-m)x)
\end{aligned}
\]
gilt. Daraus ist ersichtlich, dass $X''$ aus $X'$ hervorgeht, wenn man für jeden Term $\delta^\lambda\sin\lambda x$
\[
\begin{aligned}
&-\operatorname{cotg}\mu a\,\delta^\lambda\cos\lambda x+\operatorname{cosec}\mu a\,\delta^\lambda\cos(\mu-\lambda)x
\end{aligned}
\]
oder
\[
\begin{aligned}
&\frac{-\delta^\lambda\cos\mu a\cos\lambda x+\delta^\lambda\cos(\mu-\lambda)x}{\sin\mu a}
\end{aligned}
\]
setzt. Den Vergleich dieser Regel mit den Erläuterungen in Artikel 24 überlassen wir den Lesern.

\Needspace{28\baselineskip}
\InterpolationArticle{32}
Wenn die Funktion $X$ mit $\cos nx$ nicht abbricht, sondern weiter fortläuft, wobei die folgenden Terme durch $\alpha^\mu\cos\mu x+\alpha^{\mu+1}\cos(\mu+1)x+\text{usw.}$ ausgedrückt werden, werden wir erhalten
\SourceLeaf{323}
\[
\begin{aligned}
&\varepsilon=\alpha+\alpha^\mu\cos\mu a+\alpha^{2\mu}\cos2\mu a+\text{usw.}\\[5pt]
&\varepsilon'=\frac{1}{\sin\mu a}\{\alpha'\sin\mu a+\alpha^{\mu+1}\sin2\mu a+\alpha^{2\mu+1}\sin3\mu a+\text{usw.}\\[5pt]
&\qquad-\alpha^{2\mu-1}\sin\mu a-\alpha^{3\mu-1}\sin2\mu a-\alpha^{4\mu-1}\sin3\mu a-\text{usw.}\}\\[5pt]
&\varepsilon''=\frac{1}{\sin\mu a}\{\alpha''\sin\mu a+\alpha^{\mu+2}\sin2\mu a+\alpha^{2\mu+2}\sin3\mu+\text{usw.}\\[5pt]
&\qquad-\alpha^{2\mu-2}\sin\mu a-\alpha^{3\mu-2}\sin2\mu a-\alpha^{4\mu-2}\sin3\mu a-\text{usw.}\}\\[5pt]
&\text{usw.}\\[5pt]
&\varepsilon^{\rho-1}=\frac{1}{\sin\mu a}\{\alpha^{\mu-1}\sin\mu a+\alpha^{2\mu-1}\sin2\mu a+\alpha^{3\mu-1}\sin3\mu a+\text{usw.}\\[5pt]
&\qquad-\alpha^{\mu+1}\sin\mu a-\alpha^{2\mu+1}\sin2\mu a-\alpha^{3\mu+1}\sin3\mu a-\text{usw.}\}
\end{aligned}
\]
\Needspace{46\baselineskip}
Aus diesen Formeln lässt sich beurteilen, inwieweit der Unterschied der Koeffizienten $\varepsilon,\varepsilon',\varepsilon''$ usw. von den wahren $\alpha,\alpha',\alpha''$ usw. vernachlässigt werden darf. In diesem Fall wird also die Funktion $X''$ nicht $X$ selbst sein, mit der sie jedoch darin übereinstimmt, dass sie alle $\mu$ gegebenen Werte erfüllt. Alle ähnlichen Funktionen, die diesen Werten genügen, werden aber in der Form $X''+Y(\cos\mu x-\cos\mu a)$ enthalten sein, wobei $Y$ eine unbestimmte Funktion des Bogens $x$ von derselben Form wie $X$ bezeichnet, also ohne Sinus. Die Funktion $X''$ wird die einzige von der $\mu-1^{\mathrm{ten}}$ Ordnung sein: Jede andere ähnliche Funktion, die den gegebenen Werten genügt, wird zu einer höheren Ordnung aufsteigen. Wäre eine solche Funktion $X'''$ anderweitig bekannt, so wären alle ähnlichen Funktionen, durch die die gegebenen Werte dargestellt werden, auch in dieser Form $X'''+Y(\cos\mu x-\cos\mu a)$ enthalten, und $Y$ könnte so bestimmt werden, dass eine auf die $\mu-1^{\mathrm{te}}$ Ordnung erniedrigte Funktion hervorgeht, die eben $X''$ sein wird. Die einfachste Methode hierzu scheint jedoch darin zu bestehen, zunächst aus $X'''$ nach Artikel 24 die Funktion $X'$ und hieraus nach Artikel 31 $X''$ abzuleiten. Es sei also irgendein Term in $X'''$ gleich $L\cos\lambda x$, und man setze $\lambda=k\mu+\lambda'$, sodass $\lambda'<\mu$ und $k$ ganzzahlig ist. Dann muss nach Artikel 24 für jenen Term in $X'$
\[
\begin{aligned}
&L\cos k\mu a\cos\lambda'x-L\sin k\mu a\sin\lambda'x,
\\[5pt]
&\qquad\text{wenn }\lambda'<\tfrac12\mu,\quad\text{jedoch nicht }\lambda'=0\\[5pt]
&L\cos(k+1)\mu a\cos(\mu-\lambda')x+L\sin(k+1)\mu a\sin(\mu-\lambda')x.
\\[5pt]
&\qquad\text{wenn }\lambda'>\tfrac12\mu\\[5pt]
&L\cos\tfrac12\mu a\cos(k+\tfrac12)\mu a\cos\tfrac12\mu x\\[5pt]
&\qquad{}+L\sin\tfrac12\mu a\cos(k+\tfrac12)\mu a\sin\tfrac12\mu x,
\\[5pt]
&\qquad\text{wenn }\lambda'=\tfrac12\mu
\end{aligned}
\]
geschrieben werden. \Needspace{38\baselineskip}
Hieraus ergibt sich in $X''$ nach dem vorigen Artikel im ersten Fall
\[
\begin{aligned}
&L\cos k\mu a\cos\lambda'x+L\sin k\mu a\operatorname{cotang}\mu a\cos\lambda'x\\[5pt]
&\qquad{}-L\sin k\mu a\operatorname{cosec}\mu a\cos(\mu-\lambda')x
\end{aligned}
\]
im zweiten Fall
\[
\begin{aligned}
&L\cos(k+1)\mu a\cos(\mu-\lambda')x-L\sin(k+1)\mu a\operatorname{cotg}\mu a\cos(\mu-\lambda')x\\[5pt]
&\qquad+L\sin(k+1)\mu a\operatorname{cosec}\mu a\cos\lambda'x
\end{aligned}
\]
\SourceLeaf{324}
im dritten Fall
\[
\begin{aligned}
&L\cos\tfrac12\mu a\cos(k+\tfrac12)\mu a\cos\tfrac12\mu x\\[5pt]
&\qquad{}-L\sin\tfrac12\mu a\cos(k+\tfrac12)\mu a\operatorname{cotg}\mu a\cos\tfrac12\mu x\\[5pt]
&\qquad+L\sin\tfrac12\mu a\cos(k+\tfrac12)\mu a\operatorname{cosec}\mu a\cos\tfrac12\mu x
\end{aligned}
\]
welcher Ausdruck sich in allen drei Fällen auf
\[
\begin{aligned}
&\frac{L\sin(k+1)\mu a\cos\lambda'x-L\sin k\mu a\cos(\mu-\lambda')x}{\sin\mu a}
\end{aligned}
\]
reduziert. Wenn aber $\lambda'=0$ ist, haben wir in $X'$ lediglich $L\cos k\mu a$, welcher Term unverändert in $X''$ beibehalten wird.

\endgroup
